
\chapter*{Appendices}
\markboth{Appendices}{}
\addcontentsline{toc}{chapter}{Appendices}

\setcounter{section}{0}
\renewcommand{\thesection}{\Alph{section}}

\section{Source Code / GitHub Link}

The complete source code, final report PDF, diagrams, and relevant resources of the \textbf{Automated Car Parking System} will be uploaded to a GitHub repository for documentation and replication purposes.

\textbf{GitHub Repository Link:} \\
\url{https://github.com/Mamim57/Automated-Car-Parking-System}

The main Arduino source code file is:

\begin{itemize}
    \item \texttt{Robo\_Car\_Parking.ino}
\end{itemize}

The source code includes IR sensor reading, RFID authentication, LCD update, servo motor-based robotic arm gate control, parking slot counting, and ESP-01S Wi-Fi based data transmission.

For security reasons, sensitive information such as Wi-Fi SSID, Wi-Fi password, API keys, and database credentials should not be published directly. These values should be replaced with placeholders before public submission.

\begin{lstlisting}
String ssid = "YOUR_WIFI_NAME";
String password = "YOUR_WIFI_PASSWORD";
\end{lstlisting}

\section{Team Responsibilities}

\begin{table}[H]
\centering
\small
\renewcommand{\arraystretch}{1.15}
\caption{Team Member Responsibilities}
\label{tab:team_responsibilities}
\begin{tabular}{|p{4cm}|p{9cm}|}
\hline
\textbf{Team Member} & \textbf{Contribution} \\
\hline
Md. Mizanur Rahman & Arduino coding, system logic, Firebase/online database setup, Wi-Fi data communication, and software debugging. \\
\hline
Md. Shahid Al Mamim & Hardware setup, component selection, circuit connection, robotic arm gate setup, system testing, and report writing. \\
\hline
Lamia Akter Jesmin & Survey preparation, system analysis, requirement analysis, report writing, and Wi-Fi module integration support. \\
\hline
\end{tabular}
\normalsize
\end{table}

\section{Ethics and Compliance Statements}

The project was developed for academic and prototype demonstration purposes. The survey data was used only for academic research, and personal information was not shared in the report. RFID and parking event data were used only for system testing.

The following ethical issues were considered:

\begin{itemize}
    \item Testing results were reported honestly.
    \item Related works and technical documents were cited properly.
    \item Wi-Fi passwords, API keys, and database credentials were kept private.
    \item The robotic arm gate was tested safely in a controlled lab environment.
    \item RFID and parking event data were used only for project demonstration.
\end{itemize}

In a real deployment, proper data privacy, secure communication, and user consent should be maintained.

\section{Survey 1: Initial User Need Analysis}

The first survey was conducted to understand user perception about the product concept and its possible usefulness. During the initial concept discussion, both automated car parking and an expression bot idea were mentioned. However, the final implemented project focused only on the \textbf{Automated Car Parking System}. Therefore, the expression bot was treated as an early brainstorming idea and was not included in the final implementation.

\subsection*{Survey Questions}

\begin{enumerate}
    \item What is your job role?
    \item Are you involved in purchasing or decision making?
    \item Have you understood the product concept?
    \item Do you think the product will be helpful for your workplace?
    \item Do you have any additional comments?
\end{enumerate}

\subsection*{Survey Response Summary}

\begin{table}[H]
\centering
\small
\renewcommand{\arraystretch}{1.15}
\caption{Summary of Survey 1 Response}
\label{tab:survey1_summary}
\begin{tabular}{|p{5cm}|p{8cm}|}
\hline
\textbf{Question / Area} & \textbf{Response Summary} \\
\hline
Respondent Role & Assistant Coordinator in Web Service and Front-end Software Development. \\
\hline
Decision Making Involvement & Yes. \\
\hline
Understanding of Concept & Rated 3 out of 5. \\
\hline
Workplace Helpfulness & Rated 3 out of 5. \\
\hline
Additional Comment & The respondent found the concept interesting and mentioned that automated car parking has practical value if implemented properly. The main challenge identified was system reliability and real-time response. \\
\hline
\end{tabular}
\normalsize
\end{table}

\subsection*{Interpretation of Survey 1}

The first survey showed that the automated parking idea has practical value, but the concept needed clearer explanation and reliable implementation. Since the expression bot idea was not implemented, it was excluded from the final project scope. The feedback helped the team focus on system stability, real-time response, and proper hardware-software integration for the Automated Car Parking System.

\section{Survey 2: Prototype Feedback}

The second survey was conducted after presenting the Automated Car Parking System concept to a respondent from a garage-related professional background. The purpose was to evaluate product understanding, usefulness, purchase interest, price acceptability, and possible improvement areas.

\subsection*{Survey Questions}

\begin{enumerate}
    \item What is your job role?
    \item Have you understood the product concept?
    \item Do you think the product will be helpful for your workplace?
    \item Would you purchase this product?
    \item Do you think the product price is acceptable?
    \item How much money would you spend for this type of product?
    \item Which part of this product needs improvement?
    \item Do you have any additional comments?
\end{enumerate}

\subsection*{Survey Response Summary}

\begin{table}[H]
\centering
\small
\renewcommand{\arraystretch}{1.15}
\caption{Summary of Survey 2 Response}
\label{tab:survey2_summary}
\begin{tabular}{|p{5cm}|p{8cm}|}
\hline
\textbf{Question / Area} & \textbf{Response Summary} \\
\hline
Respondent Role & Garage Operations Manager. \\
\hline
Understanding of Concept & Rated 5 out of 5. \\
\hline
Workplace Helpfulness & Rated 4 out of 5. \\
\hline
Purchase Interest & Rated 4 out of 5. \\
\hline
Price Acceptability & Rated 4 out of 5. \\
\hline
Expected Budget & The respondent did not provide a clear budget estimate. \\
\hline
Improvement Suggestion & The system should indicate which numbered parking slot the car should go to. \\
\hline
Additional Comment & The respondent stated that the system is good and that modern technology is needed to provide better customer service. \\
\hline
\end{tabular}
\normalsize
\end{table}

\subsection*{Interpretation of Survey 2}

The second survey showed more positive feedback about the Automated Car Parking System. The respondent understood the product concept clearly and considered it useful for a garage or parking-related workplace. The purchase interest and price acceptability were also positive. The main improvement suggestion was to add individual slot guidance so that the driver can know the exact parking slot number. This feature was not implemented in the current prototype but can be added in future work using individual slot sensors and a display or mobile application.

\section{Evidences of CEP}

\begin{itemize}
    \item \textbf{P1: Depth of Knowledge} -- Required knowledge of embedded systems, robotics, sensors, RFID, servo control, and IoT communication.
    \item \textbf{P2: Conflicting Requirements} -- Balanced low cost, reliable operation, smooth robotic arm movement, real-time display, and online data transmission.
    \item \textbf{P3: Depth of Analysis} -- Included analysis of sensor behavior, RFID validation, slot counting, Wi-Fi data transmission, and robotic arm movement.
    \item \textbf{P4: Familiarity of Issues} -- Parking automation is a known problem, but using a two-link robotic arm gate created project-specific challenges.
    \item \textbf{P5: Applicable Standards} -- Followed common Arduino hardware interfacing practices and used standard libraries.
    \item \textbf{P6: Stakeholder Needs} -- Considered drivers, parking managers, and system maintainers.
    \item \textbf{P7: Interdependence} -- Sensors, RFID, servo motor, LCD, Wi-Fi module, and online server were interconnected.
\end{itemize}

\section{Evidences of CEA}

\begin{itemize}
    \item \textbf{A1: Range of Resources} -- Used Arduino UNO, IR sensors, RFID module, ESP-01S, LCD, servo motor, robotic arm materials, Arduino IDE, and online database service.
    \item \textbf{A2: Level of Interaction} -- Multiple hardware and software modules interacted to complete entry and exit operations.
    \item \textbf{A3: Innovation} -- Used a two-link robotic arm gate mechanism instead of a simple fixed barrier gate.
    \item \textbf{A4: Consequences for Society and Environment} -- Can reduce waiting time, improve parking management, and support efficient use of parking spaces.
    \item \textbf{A5: Familiarity} -- Used familiar technologies but combined them in a practical robotics and IoT-based parking automation system.
\end{itemize}

